\section{Conclusion}
\label{sec:conclusion}

This paper presents \method{}, a statistical feature-transfer framework for hierarchical split federated learning. MiSFP combines support-weighted diagonal Gaussian mixtures, output-space re-estimation after the nonlinear edge segment, and deterministic class-preserving compression under component caps and optional boundary-specific byte budgets. The analysis establishes exact mass and diagonal-moment preservation, a per-merge transport bound, packet and computational costs, and conditional fixed-predictor loss bounds. Evaluation covers four classification datasets and ISIC-2018 segmentation, with matched CIFAR-100 ablations supporting the utility of component retention. Under Dir($\alpha=0.1$), MiSFP-K4 improves accuracy over H-SFP by 3.97 percentage points, with approximately 5.2\% more total communication and 92\% more prototype traffic. The low absolute accuracy under severe heterogeneity, metric-dependent fidelity results, and sensitivity to encoder extraction limit broader deployment claims. These findings support evaluating predictive utility together with representation quality, coordinate compatibility, and both prototype and synchronization costs. Future work should investigate utility-aware component allocation, coordinate alignment, controlled overhead measurements, and privacy protection for communicated summaries.